\section{Независимые методы и дальнейшие расчёты}
\label{sec:outlook}

\subsection{Какие методы применимы к сплаву}

Для сплава замещения абсолютная свободная энергия должна учитывать и движение
атомов, и перестановки Ti и Al. При фиксированных составе и форме ячейки
статистическая сумма складывается по химическим конфигурациям; каждая
конфигурация входит со своим позиционным интегралом. Кристалл Эйнштейна в
методе Френкеля--Ладда даёт аналитическую опору для позиционной части.
Если атомы навсегда закреплены за выбранной химической конфигурацией, такой
расчёт определяет её свободную энергию. Для равновесного сплава дополнительно
требуется взвесить конфигурации со всеми разрешёнными антиструктурными
дефектами. Примеры совместного учёта химических и позиционных степеней
свободы приведены в работах Ли и Скандоло для неидеальных растворов и
Менона и соавторов для Au--Cu~\cite{frenkel1984,li2024,menon2026}.

Для независимой проверки Ti--Al подходит следующий обратимый путь.
В начале все атомы взаимодействуют как Ti, но заданное число \(k\) узлов
несёт метку Al. В конце каждая метка задаёт настоящий тип атома в MTP.
При симметричных обменах меток, одинаковых фиктивных массах, фиксированных
\(T,V\) и общем ограничении на смещения разность свободных энергий равна
\begin{align}
 U_\lambda &=(1-\lambda)U_{\rm Ti}+\lambda U_{\rm MTP}, \notag\\
 F_{\rm MTP}(k)-F_{\rm Ti}
   &=-k_{\rm B}T\ln {N\choose k}
      +\int_0^1 \big\langle U_{\rm MTP}-U_{\rm Ti}
                        \big\rangle_{\lambda,k}\,d\lambda .
 \label{eq:alloy-alchemical}
\end{align}
Все энергии здесь относятся к целой ячейке. Биномиальный множитель следует
из числа равновероятных размещений меток в начальном состоянии. Если уже в
начальной точке введено дополнительное поле, выбирающее один вариант
упорядочения, для него надо вычислить собственную сумму по заполнениям и
учесть работу снятия поля. Перестановки и смещения при каждом \(\lambda\)
принимают с энергией \(U_\lambda\). Устойчивость выбранного
ГПУ-эталона на пути проверяют по структуре и по согласованности прямого
и обратного проходов. Для свободной ячейки можно
интегрировать при постоянных \(T,p\) и получить разность энергий Гиббса.
Геометрические ограничения и обработку глобального переноса следует
согласовать между началом и концом пути. Применение такого пути к Au--Cu
описано в препринте~\cite{menon2026}.

Есть две альтернативы для других постановок. В сплаве Fe--Cr Садиг и Эрхарт
ограничивали флуктуации состава в полубольшом каноническом ансамбле и
интегрировали химический потенциал. Это позволяет измерять свободную энергию
в том числе у ячеек, содержащих межфазную границу. Для получения объёмной
бинодали становятся существенными площадь границы и энергия когерентных
напряжений~\cite{sadigh2012}. Если требуется искать также конкурирующие
кристаллические структуры, можно применить вложенную выборку
(\emph{nested sampling}). Розенброк и соавторы построили таким способом
диаграмму Ag--Pd с MTP. Для применения к Ti--Al понадобились бы проверка
потенциала на всех затронутых структурах и отдельное исследование
сходимости структурного поиска~\cite{rosenbrock2021}.

Существует опубликованное сравнение на самом Ti--Al. Николс и соавторы
привели для MTP при 1000~K значения 15,2 и 23,4~ат.\% Al
для направлений \(\alpha\to\alpha_2\) и \(\alpha_2\to\alpha\).
Авторы брали последний наблюдаемый состав перед скачком внутри петли
гистерезиса в ячейке из 1024 атомов. На эти значения влияют время
наблюдения и зарождение новой фазы. В настоящей работе поле определено
по равенству статистических весов на один вариант упорядочения.
Для сопоставления работ следует сохранить оба критерия и указать метод
рядом с каждым числом~\cite{nichols2026}.

\subsection{Последовательность контрольных экспериментов}

\paragraph{1. Разложить расхождение с C46 при 1000~K.}
Начать с одинаковых \(p=0\) и фиксированного \(c/a=1{,}56374\).
В этой геометрии для границы \(\alpha\) остаются примерно \(1{,}12\)
процентного пункта, тогда как эффект освобождения \(c/a\) уже измерен.
Из исходных фононных частот тех же обучающих структур и объёмов получить
два набора меток. Для устойчивой моды частоты \(\omega\) они равны
\[
 f_{\rm cl}=k_{\rm B}T\ln(\beta\hbar\omega),\qquad
 f_{\rm q}=k_{\rm B}T\ln\!\left[2\sinh(\beta\hbar\omega/2)\right],
 \qquad \beta=(k_{\rm B}T)^{-1}.
\]
Сохранить набор кластеров, веса регрессии, обработку нулевых мод и
объёмную сетку. Повторное определение обеих границ с классическими
метками покажет, как замена статистики ядер меняет предсказание C46 при
неизменной схеме подгонки. Затем на независимых декорациях из обеих фаз
сопоставить предсказание классического CE с явно вычисленным гармоническим
вкладом того же MTP. В малой ячейке можно выполнить полную сумму по
заполнениям. Для большей ячейки коррекцию энергии по декорациям можно
перевзвесить при достаточном перекрытии выборок. Разность свободных
энергий, её производная по составу и сдвиг обеих границ дадут оценку
погрешности кластерной аппроксимации.

Для оставшегося ангармонического вклада сравнить гармонический эталон
с прямым классическим MTP в \emph{одной и той же} области \(W\).
Для устойчивой декорации с гауссовским эталоном переход к этой области
добавляет \(-k_{\rm B}T\ln P_{\rm harm}(W<W_{\max})\) к её свободной
энергии. Эту поправку следует отслеживать отдельно от изменения силового
закона. Декорации с неустойчивыми модами требуют собственного
ангармонического эталона либо явного изменения области сопоставления.
В результате получим цепочку сравнений по обеим границам. Численные
вклады относятся к выбранному порядку замены моделей.

\paragraph{2. Замкнуть независимый цикл свободных энергий.}
Выбрать по нескольку составов по обе стороны от каждого контакта, провести
интегрирование по \(\lambda\) в обе стороны с обменом типов и контролем
порядка. Сопоставить разность свободных энергий между выбранными составами
с интегралом \(\int\Delta\mu\,dx\) по той же ветви и в тех же единицах
на атом. Если ветвь фиксируется дополнительным смещающим полем, учесть
его работу и при интегрировании, и при восстановлении исходных весов.
При резкой перестройке порядка внутри пути можно добавить окна по \(Q\).
Критерием согласия служит замыкание цикла в пределах суммарного интервала
ошибок по независимым траекториям. Так можно проверить общую нормировку
фаз и одновременно оценить практическую стоимость прямого TDI.

\paragraph{3. Определить поправку исходного MTP к DFT.}
Сформировать выборку термически декоррелированных конфигураций с сохранением
координат, ячеек и химических типов. Нужны состояния обеих фаз в окрестности
их равновесных составов. Для каждого состояния вычислить полную энергию
MTP и DFT на \emph{одинаковых} координатах; исходной целью DFT выбрать
функционал и энергетическую конвенцию обучающего набора потенциала
Qi~\cite{qi2023}. Начать с ячеек, доступных для DFT, и проверить
переносимость поправки по размеру, составу и форме ячейки. Для фиксированных
состава, фазы и механического ансамбля формула возмущений даёт
\[
 \delta G_i=-k_{\rm B}T
   \ln\left\langle \exp[-\beta(U_{\rm DFT}-U_{\rm MTP})]\right\rangle_{i,{\rm MTP}} .
\]
Здесь давление, область фазы, ограничения и мера одинаковы в обоих
ансамблях. При фиксированном объёме та же формула даёт \(\delta F_i\).
В экспоненте стоит разность энергий \emph{всей ячейки}. Оценить долю
веса крупнейших конфигураций, эффективный объём выборки и устойчивость
к разбиению её по независимым репликам. Если вес сосредоточен в
единичных кадрах, потребуются промежуточные гамильтонианы и
двусторонняя оценка. Подобная привязка свободных энергий к DFT была
продемонстрирована для сплава TaVCrW~\cite{zhu2024}.

После построения \(\delta g_i(x)=\delta G_i(x)/N\) следует снова найти
общую касательную. Для масштаба задачи малая \emph{разность} поправок
между фазами в 1~мэВ/атом при измеренном разрыве состава
\(\Delta x\simeq0{,}0746\) сдвигает поле сосуществования примерно на
\(0{,}001/0{,}0746\simeq0{,}0134\)~эВ на замену. Это оценка
чувствительности. Сдвиг самих составов зависит ещё от производных
\(\delta g_i(x)\) и от кривизны каждой ветви.

\paragraph{4. Расширить область проверки и сопоставить с измерениями.}
Короткодействующую часть потенциала стоит проверить отдельными DFT-точками
на траектории сближения Ti--Ti. Для варианта с исправленным отталкиванием
надо заново измерить границы и их зависимость от \(W\), сохранив результаты
исходного MTP как отдельный набор. Для температурного хода сольвуса
целесообразно сначала получить независимые опорные точки по обе стороны
от 1000~K. Между ними можно интегрировать уравнение Гиббса--Дюгема
при постоянном давлении~\cite{nam2007},
\[
 \frac{d(\beta\Delta\mu)}{d\beta}
     =\frac{h_{\alpha_2}-h_\alpha}{x_{\alpha_2}-x_\alpha}.
\]
Здесь \(h_i\) и \(\Delta\mu\) используют согласованную термодинамическую
конвенцию и гамильтониан без явной зависимости от температуры. При слиянии
составов правую часть труднее оценить, и требуется новый прямой расчёт.
С экспериментом следует сопоставлять обе ветви при совпадающих температурах
и в атомных процентах с указанием состава примесей, времени отжига и
погрешности определения фазовых составов. Работа Онумы и соавторов
содержит первичные измерения на Ti--Al с низким содержанием кислорода и
отдельную термодинамическую оценку~\cite{ohnuma2000}. После привязки к DFT
останутся вопросы выбора функционала, электронного температурного вклада
и условий эксперимента. Они задают предел интерпретации сравнения с опытом.
